%========================================
% MatinBook — Stage 07: Boxes Test
% Version: v1.1
% Purpose:
%   Validate the two box families and all 10 colored box types:
%     1. Family 1 — matinbox (Boyer-style, solid title bar)
%        Used by: theorem, lemma, corollary, proposition,
%                 definition, remark, exercise
%     2. Family 2 — matin-outline (RTL outline, right vertical
%        line, two horizontal rules)
%        Used by: example, proof, solution
%     3. All 5 color families (blue, green, orange, red, purple)
%     4. Shared counter behavior
%     5. Cross-references to all boxes
%
% Compilation:
%   xelatex -shell-escape stage07-boxes.tex
%   xelatex -shell-escape stage07-boxes.tex
%   (2 passes required for cover + cross-references)
%
% Expected outcome:
%   A professionally typeset PDF demonstrating that every
%   box type renders correctly, that the two families are
%   visually distinct, and that color families are consistent.
%========================================

%------------------------------------------------------------------------
% Search Path (for tests directory)
%------------------------------------------------------------------------
\makeatletter
\def\input@path{{../}}
\makeatother

%------------------------------------------------------------------------
% Document Class
%------------------------------------------------------------------------
\documentclass{matinbook}

%------------------------------------------------------------------------
% Book Metadata
%------------------------------------------------------------------------
\title{آزمون جعبه‌های رنگی}
\author{تیم توسعه MatinBook}
\date{مهر ۱۴۰۵}

\booktitle{آزمون جعبه‌های رنگی}

\renewcommand{\repository}{github.com/matinmahmoudi-cs/matinbook}

%========================================================================
% DOCUMENT
%========================================================================
\begin{document}

%------------------------------------------------------------------------
% FRONT COVER
%------------------------------------------------------------------------
\makecover
    {آزمون جعبه‌های رنگی}
    {ارزیابی دو خانواده‌ی جعبه و ده نوع محیط علمی}
    {تیم توسعه MatinBook}
    {مهر ۱۴۰۵}

%------------------------------------------------------------------------
% FRONT MATTER (symmetric margins)
%------------------------------------------------------------------------
\frontmatter
\frontmattergeometry

% Title page
\maketitle

% Copyright / Colophon
\thispagestyle{empty}
\vfill
\begin{center}
    {\large\textbf{شناسنامه آزمون}}

    \vspace{0.8cm}

    \begin{tabular}{rl}
        عنوان: & آزمون جعبه‌های رنگی \\
        نسخه: & \matinbookversion \\
        تاریخ: & \matinbookdate \\
        موتور: & \lr{XeLaTeX} \\
        مجوز: & \lr{MIT} \\
    \end{tabular}

    \vspace{0.8cm}

    {\small این سند بخشی از مجموعه آزمون‌های یکپارچگی \lr{MatinBook} است.}
\end{center}
\clearpage

% Preface
\chapter*{پیشگفتار}
\addcontentsline{toc}{chapter}{پیشگفتار}

این سند، هفتمین آزمون از مجموعه‌ی آزمون‌های یکپارچگی
\lr{MatinBook} نسخه‌ی \lr{\matinbookversion} است. هدف آن،
اعتبارسنجی کامل سیستم جعبه‌های رنگی کلاس در پنج حوزه‌ی
متمایز است: خانواده‌ی \lr{matinbox}، خانواده‌ی
\lr{matin-outline}، پنج خانواده‌ی رنگی، شمارنده‌ی مشترک
و ارجاع متقابل به همه‌ی جعبه‌ها.

در کتاب‌های ریاضی و فنی، تمایز بصری میان انواع محتوا
بسیار مهم است. قضیه، تعریف، مثال، نکته، تمرین، اثبات و
راه‌حل، هر یک نقش معنایی مشخصی دارند و باید در قالب
بصری متمایز نمایش داده شوند. \lr{MatinBook} با دو
خانواده‌ی جعبه و پنج خانواده‌ی رنگی، این تمایز را
ایجاد می‌کند.

هر بخش از این آزمون، با مثال‌های عملی و ارجاعات متقابل
همراه است تا خواننده بتواند درستی رفتار هر زیرسیستم را
در بافت واقعی یک کتاب ارزیابی کند.

\clearpage

% Table of Contents
\tableofcontents
\clearpage

% List of Figures
\listoffigures
\clearpage

% List of Tables
\listoftables
\clearpage

%------------------------------------------------------------------------
% MAIN MATTER (wide margin for notes)
%------------------------------------------------------------------------
\mainmatter
\mainmattergeometry

%========================================================================
% CHAPTER 1 — BLUE FAMILY (THEOREMS)
%========================================================================
\chapter{خانواده‌ی آبی: قضایا و اثبات‌ها}
\label{chap:family-blue}

\section{راهنمای رنگ‌ها}
\label{sec:family-blue-guide}

خانواده‌ی آبی برای مفاهیم اصلی و اثبات‌شده استفاده می‌شود.
این خانواده شامل چهار نوع جعبه‌ی \lr{matinbox} است که
از تیره به روشن مرتب شده‌اند:

\begin{itemize}
    \item \textcolor{matindarkblue}{\textbf{گزاره}} —
    \lr{matinblue!90!black} (تیره‌ترین)
    \item \textcolor{matinblue}{\textbf{قضیه}} —
    \lr{matinblue} (اصلی)
    \item \textcolor{matinblue!85!black}{\textbf{لم}} —
    \lr{matinblue!85!black}
    \item \textcolor{matinblue!70!black}{\textbf{نتیجه}} —
    \lr{matinblue!70!black} (روشن‌ترین)
\end{itemize}

\mnote{خانواده‌ی آبی، رنگ غالب سیستم است و برای محیط‌های نظری به کار می‌رود.}

\section{نمونه‌های خانواده‌ی آبی}
\label{sec:family-blue-samples}

\begin{theorem}[قضیه‌ی فیثاغورث]
    در مثلث قائم‌الزاویه، مربع وتر برابر است با مجموع
    مربعات دو ضلع دیگر:
    \[
        a^{2} + b^{2} = c^{2}
    \]
    \label{thm:pythagoras}
\end{theorem}

\begin{lemma}[لم کمکی]
    این یک لم با رنگ آبی متوسط است. لم‌ها معمولاً برای
    اثبات قضایای بزرگ‌تر به کار می‌روند و ارزش مستقلی
    ندارند.
    \label{lem:sample}
\end{lemma}

\begin{corollary}[نتیجه‌ی مستقیم]
    این یک نتیجه با رنگ آبی روشن است. نتایج معمولاً
    بلافاصله از قضایا یا لم‌ها نتیجه‌گیری می‌شوند.
    \label{cor:sample}
\end{corollary}

\begin{proposition}[گزاره‌ی نمونه]
    این یک گزاره با رنگ آبی تیره است. گزاره‌ها احکامی
    هستند که ارزش اثبات دارند اما لزوماً به اندازه‌ی
    قضایا مهم نیستند.
    \label{prop:sample}
\end{proposition}

طبق \cref{thm:pythagoras}، لم \cref{lem:sample}،
نتیجه‌ی \cref{cor:sample} و گزاره‌ی \cref{prop:sample}
همگی از یک شمارنده‌ی مشترک استفاده می‌کنند.

%========================================================================
% CHAPTER 2 — GREEN FAMILY (DEFINITIONS)
%========================================================================
\chapter{خانواده‌ی سبز: تعاریف و مفاهیم}
\label{chap:family-green}

\section{راهنمای رنگ}
\label{sec:family-green-guide}

خانواده‌ی سبز برای تعاریف و معرفی مفاهیم جدید استفاده
می‌شود. رنگ سبز نماد رشد و یادگیری است. تعریف‌ها از
خانواده‌ی \lr{matinbox} با رنگ \lr{matingreen!85!black}
استفاده می‌کنند.

\section{نمونه‌های تعاریف}
\label{sec:family-green-samples}

\begin{definition}[مشتق تابع]
    مشتق تابع $f$ در نقطه‌ی $x_{0}$ به صورت زیر تعریف
    می‌شود:
    \[
        f'(x_{0}) = \lim_{h \to 0} \frac{f(x_{0} + h) - f(x_{0})}{h}
    \]
    مشتق نشان‌دهنده‌ی نرخ تغییرات لحظه‌ای تابع است.
    \label{def:derivative}
\end{definition}

\begin{definition}[انتگرال معین]
    انتگرال معین تابع $f$ روی بازه‌ی $[a, b]$:
    \[
        \int_{a}^{b} f(x) \dd{x} = \lim_{n \to \infty}
        \sum_{i=1}^{n} f(x_{i}^{*}) \Delta x
    \]
    انتگرال سطح زیر نمودار تابع را محاسبه می‌کند.
    \label{def:integral}
\end{definition}

\begin{definition}[ماتریس]
    ماتریس $A$ با ابعاد $m \times n$ آرایه‌ای مستطیلی
    از اعداد با $m$ سطر و $n$ ستون است:
    \[
        A = [a_{ij}]_{m \times n}
    \]
    \label{def:matrix}
\end{definition}

تعریف‌های \cref{def:derivative}، \cref{def:integral} و
\cref{def:matrix} سه مفهوم بنیادین حسابان و جبر خطی را
معرفی می‌کنند.

%========================================================================
% CHAPTER 3 — ORANGE FAMILY (EXAMPLES)
%========================================================================
\chapter{خانواده‌ی نارنجی: مثال‌ها}
\label{chap:family-orange}

\section{راهنمای رنگ}
\label{sec:family-orange-guide}

خانواده‌ی نارنجی برای مثال‌های کاربردی استفاده می‌شود.
رنگ گرم نارنجی توجه خواننده را به کاربرد مفاهیم جلب
می‌کند. مثال‌ها از خانواده‌ی \lr{matin-outline} با رنگ
\lr{matinorange!90!black} استفاده می‌کنند: خط عمودی
سمت راست، دو خط افقی زیر عنوان و زیر محتوا، بدون
پس‌زمینه‌ی رنگی.

\section{نمونه‌های مثال}
\label{sec:family-orange-samples}

\begin{example}[مشتق چندجمله‌ای]
    مشتق تابع $f(x) = x^{3} - 2x^{2} + 5x - 1$:
    \[
        f'(x) = 3x^{2} - 4x + 5
    \]
    \label{ex:derivative-poly}
\end{example}

\begin{example}[محاسبه‌ی انتگرال]
    مساحت زیر نمودار $f(x) = x^{2}$ از $x = 0$ تا $x = 2$:
    \[
        \int_{0}^{2} x^{2} \dd{x}
        = \left[ \frac{x^{3}}{3} \right]_{0}^{2}
        = \frac{8}{3} \approx 2.67
    \]
    \label{ex:integral-x2}
\end{example}

\begin{example}[ضرب ماتریس‌ها]
    \[
    \mat{1 & 2 \\ 3 & 4} \times
    \mat{5 & 6 \\ 7 & 8} =
    \mat{1\cdot5+2\cdot7 & 1\cdot6+2\cdot8 \\
         3\cdot5+4\cdot7 & 3\cdot6+4\cdot8} =
    \mat{19 & 22 \\ 43 & 50}
    \]
    \label{ex:matrix-multiply}
\end{example}

مثال‌های \cref{ex:derivative-poly}، \cref{ex:integral-x2}
و \cref{ex:matrix-multiply} سه کاربرد عملی از مفاهیم
معرفی‌شده در فصل‌های قبل را نشان می‌دهند.

%========================================================================
% CHAPTER 4 — RED FAMILY (REMARKS)
%========================================================================
\chapter{خانواده‌ی قرمز: نکات و هشدارها}
\label{chap:family-red}

\section{راهنمای رنگ}
\label{sec:family-red-guide}

خانواده‌ی قرمز برای نکات مهم، هشدارها و توجه‌ها استفاده
می‌شود. رنگ قرمز به‌طور طبیعی جلب توجه می‌کند. نکات
از خانواده‌ی \lr{matinbox} با رنگ \lr{matinred!80!black}
استفاده می‌کنند.

\mnote{استفاده‌ی سنجیده از قرمز ضروری است؛ استفاده‌ی بیش از حد، اثر هشداردهندگی آن را از بین می‌برد.}

\section{نمونه‌های نکات}
\label{sec:family-red-samples}

\begin{remark}[نکته‌ی مهم]
    در محاسبه‌ی مشتق، توجه داشته باشید که
    $\deriv{}{x}(x^{n}) = nx^{n-1}$ فقط برای $n \neq 0$
    معتبر است. برای $n = 0$، مشتق تابع ثابت صفر است.
    \label{rem:derivative-power}
\end{remark}

\begin{remark}[هشدار]
    تقسیم بر صفر تعریف نشده است! هنگام ساده‌سازی
    عبارات کسری، همواره فرض کنید مخرج غیرصفر است.
    \label{rem:division-zero}
\end{remark}

\begin{remark}[نکته‌ی کاربردی]
    برای حفظ کردن فرمول مشتق توابع مثلثاتی، به یاد
    داشته باشید که مشتق سینوس، کسینوس و مشتق کسینوس،
    \textbf{منفی} سینوس است:
    \[
        \deriv{}{x}(\sin x) = \cos x, \qquad
        \deriv{}{x}(\cos x) = -\sin x
    \]
    \label{rem:trig-derivatives}
\end{remark}

%========================================================================
% CHAPTER 5 — PURPLE FAMILY (EXERCISES)
%========================================================================
\chapter{خانواده‌ی بنفش: تمرین‌ها}
\label{chap:family-purple}

\section{راهنمای رنگ}
\label{sec:family-purple-guide}

خانواده‌ی بنفش برای تمرین‌ها و مسائل استفاده می‌شود.
رنگ بنفش محرک خلاقیت و تفکر است. تمرین‌ها از خانواده‌ی
\lr{matinbox} با رنگ \lr{matinpurple!85!black} استفاده
می‌کنند.

\section{نمونه‌های تمرین}
\label{sec:family-purple-samples}

\begin{exercise}[مشتق توابع مثلثاتی]
    مشتق توابع زیر را محاسبه کنید:
    \begin{enumerate}
        \item $f(x) = \sin(2x)$
        \item $g(x) = \cos(x^{2})$
        \item $h(x) = \tan x$
    \end{enumerate}
    \label{exc:trig-derivatives}
\end{exercise}

\begin{solution}[راه‌حل]
    با استفاده از قاعده‌ی زنجیره‌ای:
    \begin{enumerate}
        \item $f'(x) = 2\cos(2x)$
        \item $g'(x) = -2x\sin(x^{2})$
        \item $h'(x) = \sec^{2} x = \frac{1}{\cos^{2} x}$
    \end{enumerate}
\end{solution}

\begin{exercise}[انتگرال‌گیری]
    انتگرال‌های زیر را محاسبه کنید:
    \begin{enumerate}
        \item $\int (3x^{2} - 4x + 1) \dd{x}$
        \item $\int \sin x \dd{x}$
        \item $\int e^{2x} \dd{x}$
    \end{enumerate}
    \label{exc:integrals}
\end{exercise}

\begin{solution}[راه‌حل]
    \begin{enumerate}
        \item $\int (3x^{2} - 4x + 1) \dd{x}
        = x^{3} - 2x^{2} + x + C$
        \item $\int \sin x \dd{x} = -\cos x + C$
        \item $\int e^{2x} \dd{x}
        = \frac{1}{2}e^{2x} + C$
    \end{enumerate}
\end{solution}

راه‌حل‌های \cref{exc:trig-derivatives} و
\cref{exc:integrals} از خانواده‌ی \lr{matin-outline}
با رنگ \lr{matingreen!70!black} استفاده می‌کنند.

%========================================================================
% CHAPTER 6 — ALL BOXES AT A GLANCE
%========================================================================
\chapter{نمای کلی همه‌ی جعبه‌ها}
\label{chap:all-boxes}

\section{مقایسه‌ی بصری}
\label{sec:all-boxes-comparison}

در این بخش، هر ده نوع جعبه را بدون محتوای طولانی در
کنار هم می‌بینیم تا تفاوت رنگ‌ها و خانواده‌ها مشهود باشد:

\subsection*{خانواده‌ی matinbox (نوار رنگی توپر)}

\begin{theorem}
    \textbf{قضیه (آبی)}: مهم‌ترین احکام ریاضی
\end{theorem}

\begin{lemma}
    \textbf{لم (آبی متوسط)}: احکام کمکی برای اثبات
\end{lemma}

\begin{corollary}
    \textbf{نتیجه (آبی روشن)}: نتایج مستقیم قضایا
\end{corollary}

\begin{proposition}
    \textbf{گزاره (آبی تیره)}: احکام با اهمیت کمتر
\end{proposition}

\begin{definition}
    \textbf{تعریف (سبز)}: معرفی مفاهیم جدید
\end{definition}

\begin{remark}
    \textbf{نکته (قرمز)}: توضیحات و هشدارهای مهم
\end{remark}

\begin{exercise}
    \textbf{تمرین (بنفش)}: مسائل برای حل کردن
\end{exercise}

\subsection*{خانواده‌ی matin-outline (خط عمودی سمت راست)}

\begin{example}
    \textbf{مثال (نارنجی)}: کاربرد مفاهیم
\end{example}

\begin{proof}
    \textbf{اثبات (خاکستری تیره)}: استدلال ریاضی
\end{proof}

\begin{solution}
    \textbf{راه‌حل (سبز)}: پاسخ تشریحی تمرین‌ها
\end{solution}

%========================================================================
% CHAPTER 7 — SUMMARY
%========================================================================
\chapter{خلاصه‌ی نتایج}
\label{chap:summary}

\section{وضعیت جعبه‌های رنگی}
\label{sec:summary-status}

جدول \cref{tab:boxes-status} وضعیت تمام جعبه‌های رنگی
\lr{MatinBook} را نشان می‌دهد.

\begin{matintable}{وضعیت تمام جعبه‌های رنگی \lr{MatinBook}}{tab:boxes-status}
\begin{tblr}{
    width = \textwidth,
    colspec = {l l l X[c]},
    hlines, vlines,
    colsep = 8pt,
    rowsep = 4pt,
    row{1} = {font=\bfseries},
}
خانواده & نوع جعبه & استایل & وضعیت \\
آبی & \lr{theorem} & \lr{matinbox} & \textcolor{matingreen}{فعال} \\
آبی & \lr{lemma} & \lr{matinbox} & \textcolor{matingreen}{فعال} \\
آبی & \lr{corollary} & \lr{matinbox} & \textcolor{matingreen}{فعال} \\
آبی & \lr{proposition} & \lr{matinbox} & \textcolor{matingreen}{فعال} \\
سبز & \lr{definition} & \lr{matinbox} & \textcolor{matingreen}{فعال} \\
قرمز & \lr{remark} & \lr{matinbox} & \textcolor{matingreen}{فعال} \\
بنفش & \lr{exercise} & \lr{matinbox} & \textcolor{matingreen}{فعال} \\
نارنجی & \lr{example} & \lr{matin-outline} & \textcolor{matingreen}{فعال} \\
خاکستری & \lr{proof} & \lr{matin-outline} & \textcolor{matingreen}{فعال} \\
سبز & \lr{solution} & \lr{matin-outline} & \textcolor{matingreen}{فعال} \\
\end{tblr}
\end{matintable}

\section{معیارهای ارزیابی}
\label{sec:summary-criteria}

در این آزمون، سیستم جعبه‌های رنگی \lr{MatinBook} بر
اساس پنج معیار ارزیابی شد:

\begin{enumerate}
    \item \textbf{تمایز خانوادگی}: تفاوت بصری میان دو خانواده.
    \item \textbf{انسجام رنگی}: هماهنگی رنگ‌ها در هر خانواده.
    \item \textbf{شمارنده‌ی مشترک}: شماره‌گذاری پیوسته.
    \item \textbf{شکست صفحات}: رفتار جعبه‌های بلند.
    \item \textbf{ارجاع هوشمند}: \lr{\textbackslash cref} با نام‌های فارسی.
\end{enumerate}

\section{نتیجه‌گیری}
\label{sec:summary-conclusion}

\textbf{تمام ده نوع جعبه‌ی رنگی \lr{MatinBook} نسخه‌ی
\lr{\matinbookversion} با موفقیت بارگذاری و آزمایش شدند.
پنج خانواده‌ی رنگی و دو خانواده‌ی جعبه، تمایز بصری
مناسبی ایجاد می‌کنند و همگی به‌درستی کار می‌کنند.}

%------------------------------------------------------------------------
% BACK MATTER (symmetric margins)
%------------------------------------------------------------------------
\backmatter
\frontmattergeometry

% Bibliography (empty placeholder)
\printbibliography[title={منابع و مراجع}]

% Index
\printindex

%------------------------------------------------------------------------
% BACK COVER
%------------------------------------------------------------------------
\authorbio{%
    تیم توسعه‌ی \lr{MatinBook}، گروهی از پژوهشگران و برنامه‌نویسان
    فارسی‌زبان است که با هدف ارائه‌ی یک راه‌حل حرفه‌ای برای
    حروف‌چینی کتاب‌های فنی فارسی فعالیت می‌کند. این پروژه حاصل
    سال‌ها تجربه در حوزه‌ی نشر دیجیتال و برنامه‌نویسی است.
}

\makebackcover{%
    این سند، هفتمین آزمون از مجموعه‌ی آزمون‌های یکپارچگی
    \lr{MatinBook} نسخه‌ی \lr{\matinbookversion} است.

    \vspace{0.5cm}

    \textbf{زیرسیستم‌های آزموده‌شده:}
    \begin{itemize}
        \item خانواده‌ی \lr{matinbox}
        \item خانواده‌ی \lr{matin-outline}
        \item پنج خانواده‌ی رنگی
        \item شمارنده‌ی مشترک
        \item ارجاع هوشمند
    \end{itemize}
}

\end{document}